\documentclass[10pt,a4paper]{amsart}
\usepackage{amsmath,amssymb,amsthm,graphicx}
% Retained minimal preamble for the Stacks Project chapter more-algebra.tex (More on Algebra).
% The chapter file pinned-original-source.tex carries only \input{preamble} on its first line, so no
% preamble text lives inside the chapter itself.  The Stacks Project uses one single shared preamble for
% every chapter; the snapshot stages it identically as sources/stacks-chow/preamble.tex and
% sources/stacks-etale-cohomology/preamble.tex.  The block below is copied mechanically, line for line,
% from sources/stacks-chow/preamble.tex lines 14-18 (the Xy-pic setup) and 154-217 (the native theorem
% declarations together with the mathematical macros).  Nothing was added, removed or redefined; the
% wrapper class and the amsmath/amssymb/amsthm/graphicx base packages are supplied by the assembler.
% Chapter: more-algebra.tex -- More on Algebra.
% For commutative diagrams we use Xy-pic
\usepackage[all]{xy}

% We use 2cell for 2-commutative diagrams.
\xyoption{2cell}

% Theorem environments.
%
\theoremstyle{plain}
\newtheorem{theorem}[subsection]{Theorem}
\newtheorem{proposition}[subsection]{Proposition}
\newtheorem{lemma}[subsection]{Lemma}

\theoremstyle{definition}
\newtheorem{definition}[subsection]{Definition}
\newtheorem{example}[subsection]{Example}
\newtheorem{exercise}[subsection]{Exercise}
\newtheorem{situation}[subsection]{Situation}

\theoremstyle{remark}
\newtheorem{remark}[subsection]{Remark}
\newtheorem{remarks}[subsection]{Remarks}

\numberwithin{equation}{subsection}

% Macros
%
\def\lim{\mathop{\mathrm{lim}}\nolimits}
\def\colim{\mathop{\mathrm{colim}}\nolimits}
\def\Spec{\mathop{\mathrm{Spec}}}
\def\Hom{\mathop{\mathrm{Hom}}\nolimits}
\def\Ext{\mathop{\mathrm{Ext}}\nolimits}
\def\SheafHom{\mathop{\mathcal{H}\!\mathit{om}}\nolimits}
\def\SheafExt{\mathop{\mathcal{E}\!\mathit{xt}}\nolimits}
\def\Sch{\mathit{Sch}}
\def\Mor{\mathop{\mathrm{Mor}}\nolimits}
\def\Ob{\mathop{\mathrm{Ob}}\nolimits}
\def\Sh{\mathop{\mathit{Sh}}\nolimits}
\def\NL{\mathop{N\!L}\nolimits}
\def\CH{\mathop{\mathrm{CH}}\nolimits}
\def\proetale{{pro\text{-}\acute{e}tale}}
\def\etale{{\acute{e}tale}}
\def\QCoh{\mathit{QCoh}}
\def\Ker{\mathop{\mathrm{Ker}}}
\def\Im{\mathop{\mathrm{Im}}}
\def\Coker{\mathop{\mathrm{Coker}}}
\def\Coim{\mathop{\mathrm{Coim}}}

% Boxtimes
%
\DeclareMathSymbol{\boxtimes}{\mathbin}{AMSa}{"02}

%
% Macros for moduli stacks/spaces
%
\def\QCohstack{\mathcal{QC}\!\mathit{oh}}
\def\Cohstack{\mathcal{C}\!\mathit{oh}}
\def\Spacesstack{\mathcal{S}\!\mathit{paces}}
\def\Quotfunctor{\mathrm{Quot}}
\def\Hilbfunctor{\mathrm{Hilb}}
\def\Curvesstack{\mathcal{C}\!\mathit{urves}}
\def\Polarizedstack{\mathcal{P}\!\mathit{olarized}}
\def\Complexesstack{\mathcal{C}\!\mathit{omplexes}}
% \Pic is the operator that assigns to X its picard group, usage \Pic(X)
% \Picardstack_{X/B} denotes the Picard stack of X over B
% \Picardfunctor_{X/B} denotes the Picard functor of X over B
\def\Pic{\mathop{\mathrm{Pic}}\nolimits}
\def\Picardstack{\mathcal{P}\!\mathit{ic}}
\def\Picardfunctor{\mathrm{Pic}}
\def\Deformationcategory{\mathcal{D}\!\mathit{ef}}
\title{Enlarging two bounded subcomplexes to a common pair that kills a derived-tensor kernel and catches an image, inside a cardinal bound}
\author{The Stacks Project Authors}
\date{stacks-project:more-algebra.tex:local-snapshot-20261012; GNU FDL 1.2 or later}
\begin{document}
\maketitle
\setcounter{subsection}{1}
\noindent\textit{Editorial note.} All delivered mathematical text is the author's own, copied from the pinned Stacks Project chapter file more-algebra.tex; no statement or proof step was rewritten, reordered or removed. The lemma carries no title of its own, so the statement span is exactly the labeled lemma environment at source lines 29894-29921 and the proof span is exactly the proof environment at source lines 29923-30009, which immediately follows it and is its own proof.\par\medskip
\noindent\textit{Editorial note.} The proof contains two occurrences of the phrase Details omitted, at source lines 29953 and 29963, concerning the size estimate for the smallest subcomplexes generated by finitely many elements and the existence of a free resolution of B of cardinality at most kappa. Both are elementary supporting steps, neither is a whole block of the main line, and the main argument from the choice of kappa to the cardinal bound of property (3) is fully present in the counted span; the phrase is preserved verbatim as written by the author.\par\medskip
\noindent\textit{Editorial note.} The chapter file pinned-original-source.tex carries only an input of the shared preamble on its first line, so no preamble text lives inside the chapter. The retained file sources/stacks-more-algebra/retained-preamble.tex is a mechanical, line-for-line copy of the project's single shared preamble as staged in the snapshot (Xy-pic setup plus the native theorem declarations and mathematical macros, including Ker, Im, colim, Tot and the bold L of the derived tensor product). Nothing was added, removed or redefined, and the excerpt serves only as preamble, never as a span coordinate.\par\medskip
\noindent\textit{Editorial note.} The source URL is the source-level Stacks Project URL, as required when more than one card is produced from one source; the per-claim tag is carried by claim\_key. The counted statement and proof refer to no labelled item of the chapter and cite no bibliography entry, so this card contributes no external reference stub and needs no bibliography.\par\medskip
\noindent\textit{Editorial note.} Property (2) as printed by the author ends in H\textasciicircum{}2(N\_2) where the surrounding text suggests H\textasciicircum{}i; the printed text is kept verbatim, since the lane copies the author's statement and neither corrects nor comments on its mathematics.\par\medskip
\noindent\textit{Allowed standard prerequisites (editorial).}\begin{itemize}
\item Arithmetic of cardinals: the maximum of a finite family of cardinals bounds the cardinality of a union or of a finite product of sets of those cardinalities. invoked at proof lines 29924-29932, where kappa = max(|A|, |B|, aleph\_0) is fixed and the statement is announced to be used without further mention. Supplies the size estimate for the smallest subcomplexes generated by finitely many elements, for the free resolution of B, and finally for the index set of the two sums that define M\_2 and N\_2, which is what makes property (3) true.
\item Every module admits a resolution by free modules of bounded cardinality, and the derived tensor product is computed by the total complex of the tensor product with such a resolution. recorded at proof lines 29962-29965 as a resolution ... -> F\textasciicircum{}\{-1\} -> F\textasciicircum{}0 -> B by free A-modules F\_i with |F\_i| at most kappa, with the cohomology of M\_1 tensor B computed by Tot(M\_1 tensor F). Turns the abstract derived tensor product into an explicit complex of cardinality at most max(kappa, |M\_1|), which is the estimate used for the cohomology of the small subcomplexes.
\item Derived tensor product commutes with filtered colimits, so a cohomology class that dies after tensoring with B already dies after tensoring with a small enough subcomplex. used at proof lines 29955-29959, where M and N are exhibited as filtered colimits of the small subcomplexes, and at proof lines 29975-29978. Produces the finite collections E\_xi and F\_xi for a kernel class xi, and the analogous collections E\_eta and F\_eta for an image class eta, inside the small colimit system.
\item The hypothesis that a : M -> N induces an isomorphism M tensor\textasciicircum{}L\_A B -> N in D(B) is an isomorphism on cohomology, and it factors through the map to N\_1. used at proof lines 29972-29974 for the kernel class xi and at proof lines 29981-29983 for the image class eta. Identifies the target of the vanishing condition for xi and the source of the preimage for eta with objects living in the N\_1 and M\_1 systems, which is what lets the enlargement be driven by finitely many elements.
\end{itemize}\medskip
\begin{lemma}
\label{lemma-enlarge}
Let $A \to B$ be a ring map. There exists a cardinal $\kappa = \kappa(A \to B)$
with the following property: Let $M^\bullet$, resp.\ $N^\bullet$
be a complex of $A$-modules, resp.\ $B$-modules. Let
$a : M^\bullet \to N^\bullet$ be a map of complexes of $A$-modules
which induces an isomorphism
$M^\bullet \otimes_A^\mathbf{L} B \to N^\bullet$ in $D(B)$.
Let $M_1^\bullet \subset M^\bullet$, resp.\ $N_1^\bullet \subset N^\bullet$
be a subcomplex of $A$-modules, resp.\ $B$-modules such that
$a(M_1^\bullet) \subset N_1^\bullet$. Then there exist
subcomplexes
$$
M_1^\bullet \subset M_2^\bullet \subset M^\bullet
\quad\text{and}\quad
N_1^\bullet \subset N_2^\bullet \subset N^\bullet
$$
such that $a(M_2^\bullet) \subset N_2^\bullet$
with the following properties:
\begin{enumerate}
\item $\Ker(H^i(M_1^\bullet \otimes_A^\mathbf{L} B) \to H^i(N_1^\bullet))$
maps to zero in $H^i(M_2^\bullet \otimes_A^\mathbf{L} B)$,
\item $\Im(H^i(N_1^\bullet) \to H^i(N_2^\bullet))$ is contained in
$\Im(H^i(M_2^\bullet \otimes_A^\mathbf{L} B) \to H^2(N_2^\bullet))$,
\item $|\bigcup M_2^i \cup \bigcup N_2^i| \leq
\max(\kappa, |\bigcup M_1^i \cup \bigcup N_1^i|)$.
\end{enumerate}
\end{lemma}

\begin{proof}
Let $\kappa = \max(|A|, |B|, \aleph_0)$. Set $|M^\bullet| = |\bigcup M^i|$
and similarly for other complexes. With this notation we have
$$
\max(\kappa, |\bigcup M_1^i \cup \bigcup N_1^i|) =
\max(\kappa, |M_1^\bullet|, |M_2^\bullet|)
$$
for the quantity used in the statement of the lemma. We are going to
use this and other observations coming from arithmetic of cardinals
without further mention.

\medskip\noindent
First, let us show that there are plenty of ``small'' subcomplexes.
For every pair of collections $E = \{E^i\}$ and $F = \{F^i\}$
of finite subsets $E^i \subset M^i$, $i \in \mathbf{Z}$
and $F^i \subset N^i$, $i \in \mathbf{Z}$ we can let
$$
M_1^\bullet \subset M_1(E, F)^\bullet \subset M^\bullet
\quad\text{and}\quad
N_1^\bullet \subset N_1(E, F)^\bullet \subset N^\bullet
$$
be the smallest subcomplexes of $A$ and $B$-modules
such that $a(M_1(E, F)^\bullet) \subset N_1(E, F)^\bullet$
and such that $E^i \subset M_1(E, F)^i$ and $F^i \subset M_2(E, F)^i$.
Then it is easy to see that
$$
|M_1(E, F)^\bullet| \leq \max(\kappa, |M_1^\bullet|)
\quad\text{and}\quad
|M_2(E, F)^\bullet| \leq \max(\kappa, |M_2^\bullet|)
$$
Details omitted. It is clear that we have
$$
M^\bullet = \colim_{(E, F)} M_1(E, F)^\bullet
\quad\text{and}\quad
N^\bullet = \colim_{(E, F)} N_1(E, F)^\bullet
$$
and the colimits are (termwise) filtered colimits.

\medskip\noindent
There exists a resolution $\ldots \to F^{-1} \to F^0 \to B$
by free $A$-modules $F_i$ with $|F_i| \leq \kappa$ (details omitted).
The cohomology modules of $M_1^\bullet \otimes_A^\mathbf{L} B$ are
computed by $\text{Tot}(M_1^\bullet \otimes_A F^\bullet)$.
It follows that
$|H^i(M_1^\bullet \otimes_A^\mathbf{L} B)| \leq \max(\kappa, |M_1^\bullet|)$.

\medskip\noindent
Let $i \in \mathbf{Z}$ and let
$\xi \in H^i(M_1^\bullet \otimes_A^\mathbf{L} B)$
be an element which maps to zero in $H^i(N_1^\bullet)$.
Then $\xi$ maps to zero in $H^i(N^\bullet)$ and hence
$\xi$ maps to zero in $H^i(M^\bullet \otimes_A^\mathbf{L} B)$.
Since derived tensor product commutes with filtered colimits,
we can find finite collections $E_\xi$ and $F_\xi$
as above such that $\xi$ maps to zero in
$H^i(M_1(E_\xi, F_\xi)^\bullet \otimes_A^\mathbf{L} B)$.

\medskip\noindent
Let $i \in \mathbf{Z}$ and let $\eta \in H^i(N_1^\bullet)$.
Then the image of $\eta$ in $H^i(N^\bullet)$ is in the image
of $H^i(M^\bullet \otimes_A^\mathbf{L} B) \to H^i(N^\bullet)$.
Hence as before, we can find finite collections $E_\eta$ and $F_\eta$
as above such that $\eta$ maps to an element of
$H^i(N_1(E_\eta, F_\eta)$ which is in the image of the map
$H^i(M_1(E_\eta, F_\eta)^\bullet \otimes_A^\mathbf{L} B) \to
H^i(N_1(E_\eta, F_\eta)$.

\medskip\noindent
Now we simply define
$$
M_2^\bullet =
\sum\nolimits_\xi M_1(E_\xi, F_\xi)^\bullet +
\sum\nolimits_\eta M_1(E_\eta, F_\eta)^\bullet
$$
where the sum is over $\xi$ and $\eta$ as in the previous two paragraphs
and the sum is taken inside $M^\bullet$. Similarly we set
$$
N_2^\bullet =
\sum\nolimits_\xi N_1(E_\xi, F_\xi)^\bullet +
\sum\nolimits_\eta N_1(E_\eta, F_\eta)^\bullet
$$
where the sum is taken inside $N^\bullet$. By construction we will
have properties (1) and (2) with these choices.
The bound (3) also follows as the set of
$\xi$ and $\eta$ has cardinality at most
$\max(\kappa, |M_1^\bullet|, |N_1^\bullet|)$.
\end{proof}

\end{document}
